\section{Related Work}
\label{sec:related}

Guo et al.~\cite{Guo2023close} asked whether ChatGPT answers can be told apart from human answers across several domains at once. They proposed HC3, a paired corpus of 85,449 question-answer rows from five English sources, with RoBERTa detectors fitted on it. The document-level detector reached 99.82 F1, which later work narrows to roughly 91.7\% under semantic-invariant tasks~\cite{Su2023plus}. One generator was collected in one window, and the paper does not ask which property of the text that score rests on.

Tian et al.~\cite{Tian2023multiscale} addressed detection when documents are too short for document-level evidence. They proposed a multiscale detector trained across text lengths and released a cleaning kit for the HC3 whitespace convention. Their appendix builds a detector out of a single test for one token identifier, which reaches 82.12 F1 at sentence level against the 81.89 they quote for a fine-tuned RoBERTa. The experiment removes one cue already known, so it bounds neither the total separability carried by surface form nor how that total compares across corpora.

Zhou et al.~\cite{Zhou2024exploiting} asked whether a transformer is necessary on the competition-era essay benchmarks. Their approach fits classical classifiers over character n-gram features instead of fine-tuning a pretrained model. On DAIGT V2 that featurisation reaches accuracy competitive with transformer detectors, which is consistent with the tie in Section~\ref{sec:results}. The finding is stated as a score rather than a decomposition, so what the n-grams read is left open. Work pairing DeBERTa with stylometric features~\cite{Yadagiri2025aigenerated} or detecting from style embeddings~\cite{Socolof2024style} likewise combines the two channels where we hold them apart.

RAID~\cite{Dugan2024raid}, M4~\cite{Wang2024m4} and SemEval-2024 Task 8~\cite{Wang2024semeval} make detector results comparable when generators, domains and attacks vary at once. Each releases labelled conditions and a leaderboard, so a detector is scored per condition rather than in aggregate. What they do not ask is what any one condition is separable by, which is the axis added here. The same blind spot is documented elsewhere, since vision datasets are identifiable from their own images~\cite{Torralba2011unbiased} and an inference model reading only the hypothesis scores far above chance~\cite{Gururangan2018annotation,Poliak2018hypothesis,Geirhos2020shortcut}. Our surface arm is that partial-input baseline moved to detection, read under the caveat of Feng et al.~\cite{Feng2019misleading}, since a high partial-input score shows a dataset is cheatable while a low one does not show it is clean.

Antoun et al.~\cite{Antoun2023towards} studied how far reported accuracy survives contact with a hostile writer. They perturbed inputs with misspellings and homoglyphs and re-scored HC3-trained detectors without retraining them. Under that attack a detector falls from 99.88 F1 to 33.57\% accuracy, and deployed detectors flag non-native English writing as machine-generated at high rates~\cite{Liang2023gpt}. Both report the fragility without locating what the detector relied on, which is what a decomposition supplies.

None of this work measures how much of a benchmark's separability is carried by surface form. Where~\cite{Tian2023multiscale} removes one cue and~\cite{Dugan2024raid,Wang2024m4} compare detectors across conditions, this paper measures each corpus with two matched arms, locates the result per sub-corpus, and checks whether the detector can read the cue at all.
